%vi: tw=78
\documentclass{article}

\usepackage{a4wide}
\usepackage[utf8x]{inputenc}
\usepackage[czech]{babel}
\usepackage[IL2]{fontenc}

\pagestyle{empty}

\begin{document}
\begin{flushright}
\today
\end{flushright}
\begin{center}
\Large\textbf{Posudek oponenta na bakalářskou práci \\
\normalsize%
\textit{Tomáš Jelínek:} Využití osobních GPS navigací při sběru a zpracování
dat}
\end{center}

Dle zadání bylo cílem práce prostudovat principy GPS, dále export dat z
GPS navigací a jejich následný import
pomocí navrhnuté a naprogramované aplikace s podporou externích mapových
podkladů do systému TripShare.

Úvod práce neuvádí do problému a nesděluje přímo, jaký je přínos práce. Je zde
pouze uvedeno, že GPS doznává v poslední době rozmachu, jak to ale souvisí s
prací, se čtenář nedozví. Poté následuje už jen popis kapitol, zcela chybí
motivace.

V další kapitole je zevrubný popis stávajících pozicovacích systémů. Rozsah
této kapitoly je dostatečný, struktura je též dobře rozvržena. Přehlednost a
orientaci v textu vylepšují vložené obrázky, nicméně ani v některých případech
není uvedeno autorství. Jsou tyto obrázky diplomantovy?

Mimo jiné je v druhé kapitole také uvedeno, že rychlost světla je neznámá
proměnná v
rovnici. Jak to autor myslel: "Další neznámé jsou posun hodin přijímače oproti
systémovému času $\Delta T$, který potřebujeme také určit, a rychlost světla
označená jako c."?

Popisu sýstémů lze vytknout snad jen, že některé ne
všeobecně známé termíny by zasluhovaly vysvětlení (např. DGPS zkratka je
použita dříve než je vysvětlena). Také u
přepočtů schází odkazy na literaturu, kam může být čtenář odkázán, chce-li
se o nich dozvědět více.

Samotné implementaci systému je v páté kapitole věnováno pouze 5 stran,
z toho část
zabírají obrázky obrazovky a příklad XML souboru. Tato kapitola by si
zasloužila rozvést. Nejméně zmínit, proč bylo použito takové rozvržení GUI,
zda se dá aplikace ovládat z příkazové řádky, proč byly zvoleny uvedené balíky
apod. Také chybí zmínka o balíku
použitém pro práci s XML, zda samotné XML je nějakým standardem TripShare a
také vysvětlení, proč obsahuje směs češtiny a angličtiny.

Aplikace samotná je psána jazykem Java, což do jisté míry (různorodost
virtuálních strojů, jiné způsob ovládání zařízení, binární JNI kód, rozdílné
GUI nebo \textit{java.lang} API např. v mobilních zařízeních, atd.) zaručuje
platformní
nezávislost. Je přehledná a uživatel si rychle zvykne na ovládací prvky.
Pravděpodobně ani úplný začátečník nebude mít obtíže aplikaci ovládat.

Poslední kapitola o převodních aplikacích je stručná, nicméně plně
dostačujicí pro účely této práce. Je zde zmíněno několik aplikací schopných
převádět jednotlivé formáty GPS dat a také komunikovat s GPS zařízeními.

Z typografického hlediska nemá práce z mého pohledu větších nedostatků.
Jedinou vadou je nemožnost kopírovat z práce text s diakritikou, pravděpodobně
bylo při sazbě využito nějakého zvláštního výstupního kódování fontu.
Práce obsahuje minimum překlepů a jeden nezvyklý termín: datumy.

Celkově práce působí propracovaně, struktuře práce byla věnována pozornost,
jednotlivé kapitoly na sebe navazují.

K samotné implementaci. Překladový systém vytvořený rozhraním
NetBeans nevytváří samostatně spustitelný \textit{.jar} balík, je-li proveden
příkaz \textit{ant jar}. Je třeba
spouštět složitějším způsobem s definicemi \textit{classpath}. Přesto se mi
program podařilo vyzkoušet a zdá se, že na testovacích datech pracuje
spolehlivě.

U obhajoby bych očekával autorovo vyjádření k následujícím bodům:
\begin{itemize}
\setlength{\itemsep}{0pt}
\item je autor tvůrcem všech obrázků?
\item proč je implementační části věnováno tak málo prostoru?
\item proč jsou elementy XML česky?
\item proč není \textit{.jar} balík samostatně spustitelný -- je to chyba
\textit{ant}u?
\end{itemize}

I přes uvedené nedostatky navrhuji práci \textbf{uznat} jako bakalářskou,
neboť autor
splnil podmínky zadání a výše uvedené výtky nejsou závažného charakteru.
Doporučuji hodnocení stupněm C.

\begin{flushright}
Jiří Slabý
\end{flushright}

\end{document}
